\QHeading{第1题　列式计算。}
\par\noindent


\Problem{① 两个数的平均数是120，其中一个数是100，另一个数是多少？}
\par


\MethodTitle{方法一}
\Answer{{\Formula{\begin{aligned}[t]&\phantom{={}}120\times2-100\\&=240-100\\&=140\end{aligned}}}\par 答：另一个数是140。}

\MethodTitle{方法二：移多补少}
\Answer{已知的100比平均数120少20，另一个数就要比120多20。}
\Answer{\Formula{\begin{aligned}[t]&\phantom{={}}120+(120-100)\\
&=120+20\\
&=140\end{aligned}}}
\Answer{答：另一个数是140。}
\Problem{② 甲、乙两数的平均数是60，丙数是30，这三个数的平均数是多少？}

\MethodTitle{方法一}
\Answer{{\Formula{\begin{aligned}[t]&\phantom{={}}(60\times2+30)\div3\\&=(120+30)\div3\\&=150\div3\\&=50\end{aligned}}}\par 答：平均数是50。}

\MethodTitle{方法二：以60为基数调整}
\Answer{先把三个数的平均数都看成60；丙实际只有30，总数少30，平均分给3个数后每个少10。}
\Answer{\Formula{\begin{aligned}[t]&\phantom{={}}60-(60-30)\div 3\\
&=60-30\div 3\\
&=60-10\\
&=50\end{aligned}}}
\Answer{答：三个数的平均数是50。}
\Problem{③ 三个连续自然数的平均数是8，这三个自然数分别是多少？}

\MethodTitle{方法一}
\Answer{中间数是8，\par $8-1=7$，$8+1=9$。\par 答：分别是7、8、9。}

\MethodTitle{方法二：由总和倒推最小数}
\Answer{三个数的总和是24。把后两个数分别减1、减2，就变成三个相同的最小数。}
\Answer{\textbf{三个数的总和}\par \Formula{\begin{aligned}[t]&\phantom{={}}8\times 3\\
&=24\end{aligned}}}
\Answer{\textbf{最小数}\par \Formula{\begin{aligned}[t]&\phantom{={}}(24-1-2)\div 3\\
&=(23-2)\div 3\\
&=21\div 3\\
&=7\end{aligned}}}
\Answer{\textbf{中间数}\par \Formula{\begin{aligned}[t]&\phantom{={}}7+1\\
&=8\end{aligned}}}
\Answer{\textbf{最大数}\par \Formula{\begin{aligned}[t]&\phantom{={}}7+2\\
&=9\end{aligned}}}
\Answer{答：这三个自然数分别是7、8、9。}
\Problem{④ 甲、乙两数之和是48，丙数是27，这三个数的平均数是多少？}

\MethodTitle{方法一}
\Answer{{\Formula{\begin{aligned}[t]&\phantom{={}}(48+27)\div3\\&=75\div3\\&=25\end{aligned}}}\par 答：平均数是25。}

\MethodTitle{方法二：先平均分两数之和，再加入第三个数}
\Answer{甲乙总和48相当于两个24。丙27比24多3，分给三个数，平均数比24多1。}
\Answer{\textbf{甲乙两数的平均数}\par \Formula{\begin{aligned}[t]&\phantom{={}}48\div 2\\
&=24\end{aligned}}}
\Answer{\textbf{三个数的平均数}\par \Formula{\begin{aligned}[t]&\phantom{={}}24+(27-24)\div 3\\
&=24+3\div 3\\
&=24+1\\
&=25\end{aligned}}}
\Answer{答：三个数的平均数是25。}
\par\noindent

\QHeading{第2题　解决问题。}

\Problem{① 下面是小夏同学星期一至星期五上学路上所花时间的统计表。}
\par

\DataTable{2}{星期 & 时间/分\\\hline
一 & 15\\\hline
二 & 17\\\hline
三 & 14\\\hline
四 & 16\\\hline
五 & 18\\\hline}
\Problem{小夏平均每天上学路上要花多长时间？}

\MethodTitle{方法一}


\MethodTitle{方法二：移多补少}
\Answer{把五天都先看成16分钟：15少1、17多1；14少2、18多2；16不增不减，正好补齐。}
\Answer{\textbf{按每天16分钟计算的总时间}\par \Formula{\begin{aligned}[t]&\phantom{={}}16\times 5\\
&=80\text{（分钟）}\end{aligned}}}
\Answer{\textbf{平均时间}\par \Formula{\begin{aligned}[t]&\phantom{={}}80\div 5\\
&=16\text{（分钟）}\end{aligned}}}
\Answer{答：平均每天上学路上要花16分钟。}
\Problem{② 琳琳参加歌唱比赛，8位评委打出的分数分别是94分、98分、95分、95分、88分、96分、93分、97分。去掉一个最高分和一个最低分，其他分数的平均数就是最后得分。琳琳的最后得分是多少？}

\MethodTitle{方法一}


\MethodTitle{方法二：以95分为基数配对}
\Answer{去掉一个98分和一个88分后，剩94、95、95、96、93、97。94与96、93与97分别平均都是95，另外两个也是95。}
\Answer{\textbf{第一对平均分}\par \Formula{\begin{aligned}[t]&\phantom{={}}(94+96)\div 2\\
&=190\div 2\\
&=95\text{（分）}\end{aligned}}}
\Answer{\textbf{第二对平均分}\par \Formula{\begin{aligned}[t]&\phantom{={}}(93+97)\div 2\\
&=190\div 2\\
&=95\text{（分）}\end{aligned}}}
\Answer{\textbf{六个分数平均分}\par \Formula{\begin{aligned}[t]&\phantom{={}}95\times 6\div 6\\
&=570\div 6\\
&=95\text{（分）}\end{aligned}}}
\Answer{答：琳琳最后得分是95分。}


\Answer{
\QHeading{计算过程：2. 解决问题}\begin{enumerate}\item 五天总用时除以5，求平均每天用时：
\Formula{\begin{aligned}[t]&\phantom{={}}(15+17+14+16+18)\div5\\&=(32+14+16+18)\div5\\&=(46+16+18)\div5\\&=(62+18)\div5\\&=80\div5\\&=16\text{（分）}\end{aligned}}
答：小夏平均每天上学路上要花16分钟。\item 最高分98分，最低分88分，去掉后剩下6个分数：94、95、95、96、93、97。\Formula{\begin{aligned}[t]&\phantom{={}}(94+95+95+96+93+97)\div6\\&=(189+95+96+93+97)\div6\\&=(284+96+93+97)\div6\\&=(380+93+97)\div6\\&=(473+97)\div6\\&=570\div6\\&=95\text{（分）}.\end{aligned}}答：琳琳的最后得分是95分。\end{enumerate}

}
